This chapter first introduces the background the thesis relies on: the irrigation method, the crop calendars, the language models and how they are adapted, the metrics, and the statistics (\cref{sec:preliminaries}). It then reviews research on agricultural language models, hallucination, abstention, clarification, Bengali language processing, test-set leakage and human evaluation (\cref{sec:review}), and closes with what that research establishes and what it leaves open (\cref{sec:lit-summary}).

\section{Preliminaries}
\label{sec:preliminaries}

\subsection{Crop water requirement and the FAO-56 method}
\label{sec:prelim-fao}

FAO Irrigation and Drainage Paper~56 gives the standard procedure for computing how much water a crop uses~\cite{allen1998fao56}. It separates the effect of weather from the effect of the crop. The weather enters through the reference evapotranspiration, $\mathrm{ET}_0$: the water that a standard, well-watered grass surface would lose to the air on a given day, in millimetres. The crop enters through a crop coefficient, $K_c$, which scales $\mathrm{ET}_0$ to the crop's own water use, $\mathrm{ET}_c = K_c \times \mathrm{ET}_0$. In the single-coefficient approach, $K_c$ changes over the season: it is low while the crop is young, highest in the middle of the season and falls again as the crop matures. FAO-56 tabulates typical values for the initial, mid-season and late-season stages of many crops, and recommends adjusting them for local climate.

A field needs irrigation to make up whatever part of $\mathrm{ET}_c$ rain has not supplied. Only part of the rain that falls is available to the crop, because some runs off, some drains below the roots and some evaporates from the surface. The available part is the effective rainfall. FAO Irrigation and Drainage Paper~25 discusses its definitions and the many methods of estimating it, from soil-water measurements to empirical rules~\cite{dastane1978effective}. Empirical rules typically apply a smaller fraction to heavy rain than to light rain. The testbed in this thesis uses such a rule with three tiers (\cref{eq:peff}). That rule, its thresholds and its fractions are the project's own convention: they are not taken from FAO-56, FAO-25 or BARC, and a model can only know them if it has seen examples of them.

\subsection{Crop calendars and irrigation in Bangladesh}
\label{sec:prelim-barc}

Dry-season irrigation in Bangladesh is supplied mostly by groundwater. Rising withdrawals have enabled dry-season rice cultivation across the Bengal Basin, and the collective pumping of millions of smallholder farmers has measurably changed how much of the monsoon the aquifers capture~\cite{shamsudduha2022bengal}. Irrigation decisions are therefore made field by field, by farmers who each choose when to run a pump and for how long.

The Bangladesh Agricultural Research Council publishes crop calendars through an online application~\cite{barc2026cropcalendar}. For each crop, a calendar gives the sowing or planting window, the growth stages and their timing, the periods when irrigation matters most, and notes on fertiliser and pests. These are fixed facts that a question-answering system can store. The inputs that change daily or from farm to farm, such as $\mathrm{ET}_0$, rainfall and the size of a particular field, are not in any calendar.

\subsection{Large language models and instruction tuning}
\label{sec:prelim-llm}

The models used in this thesis are decoder-only Transformers~\cite{vaswani2017attention}, pre-trained to predict the next token of text and then tuned to follow instructions. Instruction tuning on many tasks written as natural-language instructions improves a model's performance on tasks it was not trained on~\cite{wei2022finetuned}, and further tuning with human feedback makes models follow the intent of a user's request more closely~\cite{ouyang2022training}. The four checkpoints studied here are instruction-tuned releases of the Llama~3~\cite{grattafiori2024llama}, Mistral NeMo~\cite{mistral2024nemo}, Phi-3.5~\cite{abdin2024phi3} and Phi-4~\cite{microsoft2025phi4mini} families. Each has its own tokenizer, which splits text into subword units; the same sentence can take very different numbers of tokens in different languages (\cref{sec:rev-bengali}).

\subsection{Parameter-efficient fine-tuning}
\label{sec:prelim-peft}

Updating every weight of a model with billions of parameters needs more memory than a single consumer GPU has. Parameter-efficient methods train a small number of new parameters and leave the pre-trained weights frozen. Adapter modules insert small trainable layers into each Transformer block~\cite{houlsby2019parameter}. Low-rank adaptation (LoRA) instead adds to selected weight matrices a trainable update that is the product of two low-rank matrices, so that the update has far fewer parameters than the matrix it modifies~\cite{hu2022lora}. QLoRA goes further: the frozen base model is stored in 4-bit NormalFloat (NF4) precision and gradients flow through it into the LoRA adapters, which makes it possible to fine-tune models of this size on one GPU~\cite{dettmers2023qlora}. This thesis uses QLoRA with adapters on every projection matrix (\cref{sec:adapter-config}).

\subsection{Retrieval-augmented generation and runtime values}
\label{sec:prelim-rag}

Retrieval-augmented generation (RAG) combines a language model with a searchable document index: for each question, a retriever finds relevant passages and the model generates an answer conditioned on them~\cite{lewis2020rag}. RAG suits questions whose answers are written down somewhere, such as a crop's sowing window. It cannot supply a value that exists in no document. For a runtime-conditioned irrigation query, the day's $\mathrm{ET}_0$, the recent rainfall and the size of the farmer's field must come from the question, a sensor or a weather service. A retrieval system therefore faces the same problem as a fine-tuned one when the farmer leaves one of those values out: something has to notice that it is missing.

\subsection{Evaluation metrics}
\label{sec:prelim-metrics}

For numeric answers the natural measure is whether the stated value is correct within a tolerance, reported as the share of rows that meet it. For free-text answers there is rarely one correct wording, so overlap metrics are used. chrF scores the overlap of character $n$-grams between an answer and a reference as an F-score~\cite{popovic2015chrf}, and chrF++ adds word unigrams and bigrams~\cite{popovic2017chrfpp}. Character-level metrics suit a language such as Bengali, whose words take many inflected forms. ROUGE-L scores the longest common subsequence of words~\cite{lin2004rouge}. BERTScore compares the contextual embeddings of the tokens in an answer and a reference~\cite{zhang2020bertscore}. Word-level Dice, twice the number of shared words divided by the total number of words, is the simplest of these. Automatic overlap metrics agree only weakly with human judgements of generated text~\cite{novikova2017why}, which is why this thesis checks them against human ratings (\cref{sec:res-human}).

\subsection{Inter-rater agreement}
\label{sec:prelim-agreement}

When several people rate the same items, their agreement shows how much a rating reflects the item rather than the rater. Fleiss' $\kappa$ measures chance-corrected agreement among many raters assigning items to nominal categories~\cite{fleiss1971measuring}. Krippendorff's $\alpha$ handles any number of raters, missing ratings and different levels of measurement, including ordinal scales, and has been proposed as a standard reliability measure~\cite{hayes2007answering}. Both equal zero at chance and one at perfect agreement. Conventional bands for interpreting agreement coefficients exist~\cite{landis1977measurement}, but we report the coefficients with their confidence intervals and a threshold named before testing rather than rely on verbal labels.

\subsection{Paired significance tests and multiple comparisons}
\label{sec:prelim-stats}

When two systems answer the same test items, their results are paired and a paired test is more powerful than an unpaired one. For binary outcomes (correct or not), McNemar's test uses only the discordant items, those that one system gets right and the other wrong~\cite{mcnemar1947note}. For continuous or ordinal scores, the Wilcoxon signed-rank test ranks the paired differences~\cite{wilcoxon1945individual}. For more than two systems on the same items, the Friedman test compares ranks within each item~\cite{friedman1937use}, and Cochran's Q does the same for binary outcomes~\cite{cochran1950comparison}. The bootstrap estimates the uncertainty of any statistic by resampling the items with replacement~\cite{efron1979bootstrap}; we use it for confidence intervals of differences. Running many tests raises the chance of a false positive; Holm's procedure controls this by comparing the ordered $p$-values with thresholds that relax step by step~\cite{holm1979simple}. A survey of significance testing in natural language processing (NLP) found that it is often ignored or misused, and proposed a protocol for choosing tests~\cite{dror2018hitchhiker}.

\section{Review of Existing Research}
\label{sec:review}

\subsection{Language models for agricultural advisory}
\label{sec:rev-agri}

Several groups have studied LLMs for agricultural extension. Tzachor and colleagues assessed whether a general-purpose model could simplify scientific knowledge and give personalised, location-specific advice, tested it on technical advice for cassava farmers in Nigeria, and concluded that safe use needs a design process with human experts in the loop~\cite{tzachor2023large}. Silva and colleagues evaluated LLMs on agriculture exams from Brazil, India and the United States, and found that GPT-4 answered enough questions to pass agronomist certification exams~\cite{silva2023gpt4}. Benchmarks such as AgXQA were built specifically for agricultural extension questions~\cite{kpodo2024agxqa}.

Systems for farmers have mostly used retrieval. Farmer.Chat is a generative-AI chatbot deployed in four countries that draws on a curated knowledge base~\cite{singh2024farmerchat}, and AgroLLM compares three LLMs inside a RAG pipeline over agricultural texts~\cite{samuel2025agrollm}. Balaguer and colleagues compared RAG with fine-tuning for location-specific agricultural questions and found that each helped and that the gains added up~\cite{balaguer2024rag}. For Bangladesh, KrishokBondhu is a voice-based advisory service for Bengali-speaking farmers built on RAG over agricultural handbooks and extension manuals~\cite{krishokbondhu2025}, and a cross-lingual pipeline translates Bengali questions into English to retrieve from English manuals~\cite{hossain2026crosslingual}. KrishokChat is a Bengali agricultural benchmark built from government publications with a separate set of real farmer queries; its authors report that fine-tuning reduced refusal on safety-critical queries~\cite{krishokchat2026}.

Two features separate this literature from the present study. First, its questions are knowledge questions: what to plant, which pesticide to use, what a symptom means. The answers exist in documents. Second, its evaluations ask whether answers are correct, relevant or well grounded, not what a system does when the question lacks information the answer depends on. The KrishokChat finding that fine-tuning reduced refusal points in the same direction as the present study, but in a different setting.

\subsection{Hallucination and fine-tuning}
\label{sec:rev-hallucination}

Hallucination, text that is fluent but unsupported or wrong, was surveyed first for natural language generation in general~\cite{ji2023survey} and then for LLMs, where plausible but non-factual content is the main concern~\cite{huang2025survey}. Several studies connect it to fine-tuning. Gekhman and colleagues controlled how many fine-tuning examples introduced knowledge the model did not already have; such examples were learned more slowly than the others, and once learned, they increased the model's tendency to hallucinate roughly in proportion to their number~\cite{gekhman2024finetuning}. Kang and colleagues found that a fine-tuned model's hallucinated answers to unfamiliar questions mirror the answers attached to the unfamiliar examples in its fine-tuning data, so changing how those examples are supervised changes the behaviour~\cite{kang2025unfamiliar}. R-Tuning starts from the observation that ordinary instruction tuning makes a model complete an answer whether or not it knows it, and teaches models to decline questions outside their knowledge instead~\cite{zhang2024rtuning}.

Kalai and colleagues give a statistical account of why the behaviour persists. Training and evaluation reward guessing over admitting uncertainty, because a model graded on accuracy alone scores better by guessing, so benchmarks that ignore uncertainty keep encouraging confident answers~\cite{kalai2025why}. Kadavath and colleagues showed that large models are, when asked in the right format, reasonably well calibrated about whether they know an answer~\cite{kadavath2022language}; the problem is less that a model cannot tell when it is uncertain than that its training does not reward saying so.

These studies concern facts a model does not know. The present study concerns a different gap: a value that no model could know, because the user did not provide it. The mechanism these papers describe, that fine-tuning teaches a model to produce the kind of answer its training examples contain, is the one our results point to.

\subsection{Abstention and selective prediction}
\label{sec:rev-abstention}

A system abstains when it declines to answer. Reading-comprehension benchmarks have required abstention on questions whose answer is not in the passage~\cite{rajpurkar2018know}, and selective question answering asks a model to answer as many questions as it can while abstaining where it is likely to be wrong, including under a shift in domain~\cite{kamath2020selective}. A survey of abstention in LLMs organises the literature by the query, the model and human values~\cite{wen2024know}. Brahman and colleagues widened the scope of noncompliance beyond unsafe requests to incomplete, unsupported and indeterminate ones, found that models comply with many requests they should not, and reported that training with low-rank adapters balanced appropriate refusal against other abilities better than full fine-tuning~\cite{brahman2024art}.

Two recent benchmarks study questions that are missing information. AbstentionBench evaluates abstention across many datasets, including underspecified questions, and finds that abstention is unsolved, that fine-tuning for extended step-by-step problem solving (so-called reasoning models) degrades it, and that a carefully written system prompt helps in practice without solving the problem~\cite{kirichenko2025abstentionbench}. Fan and colleagues gave models maths questions with a missing premise; models tuned for long chains of thought wrote much longer answers and failed to abstain, while models without that tuning noticed the missing premise more often~\cite{fan2025missing}. Both studies show that a kind of fine-tuning can make models worse at recognising an unanswerable question. Neither separates abstention from clarification, and neither concerns a domain where the missing value is a routine input.

\subsection{Clarification questions and underspecified inputs}
\label{sec:rev-clarification}

Asking for clarification is the alternative to guessing and to refusing. Rao and Daum\'e ranked candidate clarification questions by how useful their answers would be~\cite{rao2018learning}, and Aliannejadi and colleagues showed that a single good clarifying question can greatly improve retrieval in conversational search~\cite{aliannejadi2019asking}. In open-domain question answering, over half of the questions in one benchmark turned out to be ambiguous~\cite{min2020ambigqa}.

With generative LLMs the question became whether models ask on their own. Kuhn and colleagues found that current models rarely ask users to clarify ambiguous questions and give incorrect answers instead; prompting a model to detect ambiguity and ask first helped~\cite{kuhn2022clam}. Zhang and Choi built a framework for deciding when a system should ask, based on the spread of plausible user intents~\cite{zhang2025clarify}. Several studies train the behaviour in. Zhang, Knox and Choi traced the problem to preference data that judges responses only by their prior context, and taught models to ask by labelling preferences with simulated future turns~\cite{zhang2024modeling}. STaR-GATE fine-tuned a model iteratively on the questions that most improved its later answers~\cite{andukuri2024stargate}. ClarifyCoder fine-tuned code models on synthetic examples in which an ambiguous programming task calls for a question before any code~\cite{wu2025clarifycoder}.

QuestBench is the closest to our clarification probe. It poses reasoning problems with exactly one missing variable and asks whether a model can pick the one question that would resolve it; models struggled to identify the right question even when they could solve the fully specified version~\cite{li2025questbench}. Our incomplete irrigation questions have the same structure: one required input is absent and one question would supply it. What we add is the effect of fine-tuning on this behaviour, a comparison with abstention inside the same training runs, and a measurement of the form the error takes.

\subsection{Bengali language processing and tokenizer cost}
\label{sec:rev-bengali}

Most of the world's languages are barely represented in language technology~\cite{joshi2020state}. Bengali is widely spoken but has long been treated as a low-resource language in NLP; BanglaBERT and its benchmark were built to address that for language understanding~\cite{bhattacharjee2022banglabert}. Evaluations of LLMs on Bengali tasks found that zero-shot models performed poorly on most tasks, and open-weight models especially so~\cite{kabir2024benllmeval}, and a recent study of open models on translated benchmarks confirmed a substantial English--Bengali gap, part of which comes from answer-format failures that differ between the languages~\cite{bhowmik2025script}.

Tokenisation adds a cost. The same text translated into different languages can take very different numbers of tokens~\cite{petrov2023language}, so speakers of some languages pay more per request, fit less text into a context window and often get poorer results~\cite{ahia2023all}. The same study of open models found that Bengali costs several times more tokens per word than English, with the factor set by the tokenizer rather than the script~\cite{bhowmik2025script}. We measure this cost for our four tokenizers (\cref{fig:tokenizer}) because it affects both training and the comparison of validation losses.

\subsection{Leakage and contamination in evaluation}
\label{sec:rev-leakage}

A test score means little if the test items, or close variants, were seen in training. Lewis, Stenetorp and Riedel found that a large share of test questions in popular open-domain QA benchmarks had near-duplicate paraphrases in the training sets, and that models did much worse on questions they could not have memorised~\cite{lewis2021question}. Pre-trained models can exploit test data present in their pre-training corpus~\cite{magar2022data}, contamination makes models look better than they are and should be measured for every benchmark~\cite{sainz2023nlp}, and methods now exist to detect it~\cite{golchin2024time}. Outside NLP, a survey of machine-learning-based science found data leakage to be widespread and a cause of severe failures to reproduce published results~\cite{kapoor2023leakage}. These studies motivated the split-first design and the four leakage detectors of this thesis, and the lesson of the second version of our own system (\cref{sec:scope}).

\subsection{Human evaluation of generated text}
\label{sec:rev-human}

Because automatic metrics agree only weakly with human judgements~\cite{novikova2017why}, generated text is often rated by people. A review of human evaluation practice found wide variation in how studies choose raters, design scales and report agreement, and proposed best-practice guidelines~\cite{vanderlee2021human}. Our human study follows several of those practices, such as reporting agreement with confidence intervals and naming the primary measures in advance, and departs from others, such as using the authors as raters, for reasons given in \cref{sec:impl-human}.

% AUTO-GENERATED by scripts/thesis/tables.py
\begin{table}[tbp]
\centering
\scriptsize
\caption{Related work compared on the dimensions this thesis studies. Entries record only what each work's abstract states; ``--'' means the abstract does not state it, not that the work lacks it. ``Missing-input probe'' means the evaluation includes questions that cannot be answered without information the user did not give.}
\label{tab:related}
\begin{tabular}{>{\raggedright\arraybackslash}p{0.2\textwidth}>{\raggedright\arraybackslash}p{0.13\textwidth}l>{\raggedright\arraybackslash}p{0.15\textwidth}>{\raggedright\arraybackslash}p{0.17\textwidth}>{\raggedright\arraybackslash}p{0.12\textwidth}}
\toprule
Work & Domain & Language & Fine-tuning studied & Missing-input probe & Human evaluation \\
\midrule
QuestBench~\cite{li2025questbench} & logic, planning, maths & en & -- & yes (one missing variable) & -- \\
AbstentionBench~\cite{kirichenko2025abstentionbench} & general QA & en & yes (reasoning-tuned vs base) & yes (underspecified) & -- \\
Missing premise~\cite{fan2025missing} & maths & en & yes (reasoning-tuned vs base) & yes (missing premise) & -- \\
Gekhman et al.~\cite{gekhman2024finetuning} & closed-book QA & en & yes & no (unknown facts) & -- \\
Kang et al.~\cite{kang2025unfamiliar} & QA, biographies & en & yes & no (unfamiliar queries) & -- \\
R-Tuning~\cite{zhang2024rtuning} & QA & en & yes & no (unknown facts) & -- \\
CLAM~\cite{kuhn2022clam} & ambiguous QA & en & no (prompting) & yes (ambiguity) & -- \\
STaR-GATE~\cite{andukuri2024stargate} & personalised tasks & en & yes & yes (unstated preferences) & -- \\
ClarifyCoder~\cite{wu2025clarifycoder} & code generation & en & yes & yes (ambiguous requirements) & -- \\
Balaguer et al.~\cite{balaguer2024rag} & agriculture & en & yes (vs RAG) & -- & -- \\
Farmer.Chat~\cite{singh2024farmerchat} & agriculture & multilingual & -- & -- & yes (deployment) \\
KrishokBondhu~\cite{krishokbondhu2025} & agriculture & bn & no (RAG) & -- & -- \\
KrishokChat~\cite{krishokchat2026} & agriculture & bn & yes (one model) & -- & -- \\
\midrule
This thesis & irrigation advisory & bn, en & yes (four models $\times$ three seeds) & yes (designed clarify probe) & yes (five annotators) \\
\bottomrule
\end{tabular}
\end{table}


\section{Summary of Key Findings}
\label{sec:lit-summary}

The literature establishes several things this thesis builds on. Fine-tuning shapes how a model behaves when it lacks information: examples that introduce new knowledge increase hallucination~\cite{gekhman2024finetuning}, models reproduce the kind of answer attached to unfamiliar training examples~\cite{kang2025unfamiliar}, and accuracy-only training and evaluation reward guessing~\cite{kalai2025why}. Models rarely ask for clarification unless they are prompted or trained to~\cite{kuhn2022clam,li2025questbench}, targeted training can teach them to ask~\cite{andukuri2024stargate,zhang2024modeling,wu2025clarifycoder}, and some fine-tuning degrades abstention on underspecified questions~\cite{kirichenko2025abstentionbench,fan2025missing}. Agricultural LLM systems, including systems for Bengali-speaking farmers, are being built and evaluated mainly for the correctness of their answers to knowledge questions~\cite{singh2024farmerchat,krishokbondhu2025,krishokchat2026}.

\Cref{tab:related} summarises how the most relevant works relate to this one. Among the works we reviewed, we did not find a study that measures, within the same fine-tuning runs, whether a behaviour present in the training data (abstention) and one absent from it (clarification) transfer; that measures whether fine-tuning changes the form of an invented answer, from visibly unreliable to authoritative, rather than only its rate; or that does either for runtime-conditioned questions in a low-resource, bilingual domain with a leakage-audited test set and a human evaluation. This thesis addresses that gap. It is an empirical extension of prior findings to a new setting, and its results should be read alongside the benchmark studies above, which cover far more models and tasks.
